\section{Modelo de comunicación entre servicios}
\emergencystretch=3em

Esta sección describe cómo se comunican los dos servicios del esqueleto que está en la carpeta \texttt{src/} del repositorio. No describe todavía el sistema completo: los nombres y puertos son provisionales hasta cerrar el ADR-002.

\subsection{Servicios del esqueleto}

\begin{table}[H]
\centering
\small
\caption{Servicios del esqueleto}
\begin{tabularx}{\linewidth}{@{}p{3.1cm}p{1.6cm}X@{}}
\toprule
\textbf{Servicio} & \textbf{Puerto} & \textbf{Responsabilidad} \\
\midrule
\texttt{convocatorias} & 8081 & Lista convocatorias de ejemplo y las publica. Al publicar una, avisa al otro servicio. \\
\texttt{notificaciones} & 8082 & Recibe avisos, los guarda en memoria y los lista. \\
\bottomrule
\end{tabularx}
\end{table}

\subsection{Mecanismo usado en esta entrega}

Los servicios se comunican por HTTP con mensajes en formato JSON. El servicio \texttt{convocatorias} actúa como cliente y \texttt{notificaciones} como servidor. Cuando se publica una convocatoria ocurre lo siguiente:

\begin{enumerate}[leftmargin=*]
  \item El cliente envía \texttt{POST /convocatorias/\{id\}/publicar} a \texttt{convocatorias}.
  \item \texttt{convocatorias} cambia el estado de la convocatoria a \emph{publicada}.
  \item \texttt{convocatorias} envía \texttt{POST /notificaciones} a \texttt{notificaciones} con el evento \texttt{ConvocatoriaPublicada}, el identificador y el título. Espera la respuesta como máximo 2 segundos.
  \item Si \texttt{notificaciones} responde, la respuesta de \texttt{convocatorias} incluye \texttt{notificacion: enviada}.
  \item Si \texttt{notificaciones} no responde, la respuesta incluye \texttt{notificacion: pendiente}. La convocatoria queda publicada y el servicio sigue atendiendo solicitudes.
\end{enumerate}

\subsection{Comportamiento observado}

Las pruebas se ejecutaron el 20 de noviembre de 2026 desde un clon limpio del repositorio, siguiendo el README de \texttt{src/}, en la máquina DESKTOP-H4C4SME. Las capturas están en la carpeta \texttt{evidencias/}.

\begin{table}[H]
\centering
\small
\caption{Resultados de las pruebas del esqueleto}
\begin{tabularx}{\linewidth}{@{}p{3.6cm}XX@{}}
\toprule
\textbf{Situación} & \textbf{Resultado observado} & \textbf{Evidencia} \\
\midrule
Ambos servicios activos & Los dos responden a \texttt{/health}. & \path|01-servicios-en-marcha.png| \\
Publicar la convocatoria 2 & Respuesta con \texttt{notificacion: enviada} y un aviso registrado en \texttt{notificaciones}. & \path|02-comunicacion-entre-servicios.png| \\
\texttt{notificaciones} detenido & Publicar la convocatoria 1 responde con \texttt{notificacion: pendiente}; \texttt{convocatorias} sigue activo. & \path|03-tolerancia-a-fallos.png| \\
\bottomrule
\end{tabularx}
\end{table}

\subsection{Interacciones del diagrama de contenedores}

La tabla recoge cada flecha del diagrama C4 de contenedores (nivel 2). \emph{Sinc.} es una llamada que espera respuesta y \emph{asinc.} deja un evento en la cola y sigue. De todo el diagrama, en el esqueleto solo está implementada la interacción de \texttt{convocatorias} con \texttt{notificaciones}, y hoy es HTTP directo; en el diagrama ese camino pasa por la cola, que es lo planeado según el ADR-002. El resto son decisiones de diseño sin implementar. Los usuarios (Vicerrectorado, Evaluadores, Postulantes y Administrador) usan el Frontend desde el navegador y no se listan como interacción entre contenedores.

{\small
\begin{longtable}{@{}p{3.7cm}p{1.3cm}p{2.9cm}p{5.6cm}@{}}
\caption{Interacciones entre contenedores y mecanismo propuesto}\\
\toprule
\textbf{Interacción} & \textbf{Tipo} & \textbf{Mecanismo} & \textbf{Justificación} \\
\midrule
\endfirsthead
\toprule
\textbf{Interacción} & \textbf{Tipo} & \textbf{Mecanismo} & \textbf{Justificación} \\
\midrule
\endhead
Frontend $\to$ API Gateway & Sinc. & HTTP con JSON & El usuario espera la respuesta en pantalla; el gateway valida el token antes de pasar la solicitud. \\
API Gateway $\to$ Postulaciones & Sinc. & HTTP con JSON & Consulta u orden que necesita respuesta inmediata, con tiempo máximo de espera. \\
API Gateway $\to$ Evaluación & Sinc. & HTTP con JSON & Igual que la anterior: el evaluador necesita ver el resultado. \\
API Gateway $\to$ Convocatorias & Sinc. & HTTP con JSON & Es el mecanismo que ya está probado en el esqueleto. \\
API Gateway $\to$ Auditoría & Sinc. & HTTP con JSON & Consulta de registros de auditoría por parte del Administrador. \\
API Gateway $\to$ Notificaciones & Sinc. & HTTP con JSON & Consulta de las notificaciones del usuario; es lectura, no aviso. \\
API Gateway $\to$ IA & Sinc. & HTTP con JSON & La persona espera el resultado del análisis; la llamada necesita tiempo máximo de espera. \\
Evaluación $\to$ Postulaciones & Sinc. & HTTP con JSON & Evaluación necesita los datos de la postulación para continuar (etiqueta \emph{consulta} del diagrama). \\
Postulaciones $\to$ Supabase Storage & Sinc. & API HTTP del servicio & Se guardan y leen los PDF de los documentos; hay que saber si la carga se completó. \\
Postulaciones $\to$ PostgreSQL & Sinc. & Conexión SQL & Cada servicio con datos tiene su propia base (ADR-001). \\
Evaluación $\to$ PostgreSQL & Sinc. & Conexión SQL & Igual: base propia por servicio. \\
Convocatorias $\to$ PostgreSQL & Sinc. & Conexión SQL & Igual: base propia por servicio. \\
Auditoría $\to$ PostgreSQL & Sinc. & Conexión SQL & Igual: base propia por servicio. \\
Convocatorias $\to$ Cola & Asinc. & Cola de mensajes (ADR-002) & Publica el evento de convocatoria publicada sin esperar a que el receptor esté disponible. \\
Evaluación $\to$ Cola & Asinc. & Cola de mensajes (ADR-002) & Publica el evento de evaluación terminada; el aviso tolera retraso. \\
Postulaciones $\to$ Cola & Asinc. & Cola de mensajes (ADR-002) & Publica el evento de postulación recibida; el aviso tolera retraso. \\
Cola $\to$ Auditoría & Asinc. & Cola de mensajes (ADR-002) & Auditoría consume los eventos cuando puede, con confirmación de recepción. \\
Cola $\to$ Notificaciones & Asinc. & Cola de mensajes (ADR-002) & Un aviso no se pierde si notificaciones está caído, que es la limitación del esqueleto. \\
Auditoría $\to$ Google Drive & Sinc. & API HTTP del servicio (respaldos) & Hay que confirmar que el respaldo se escribió. \\
Notificaciones $\to$ Correo & Sinc. & Protocolo del servicio de correo & El proveedor responde si aceptó el envío; el mecanismo exacto se define en E2. \\
IA $\to$ Proveedores de IA & Sinc. & API HTTP del proveedor & Los proveedores (Groq, Mistral, Hugging Face) se consumen por su API; hay que limitar el tiempo de espera. \\
\bottomrule
\end{longtable}
}

\subsection{Limitaciones conocidas}

\begin{itemize}[leftmargin=*]
  \item Los datos están en memoria; al reiniciar un servicio se pierden.
  \item Un aviso que queda \emph{pendiente} no se reintenta: si \texttt{notificaciones} vuelve a estar disponible, ese aviso no se envía solo.
  \item No hay autenticación entre servicios.
  \item La llamada es síncrona: \texttt{convocatorias} espera la respuesta (hasta 2 segundos) antes de contestar.
\end{itemize}

\subsection{Propuesta para el ADR-002}

La decisión definitiva del mecanismo se toma en el ADR-002 con el Arquitecto. Como punto de partida, se propone separar dos tipos de comunicación:

\begin{itemize}[leftmargin=*]
  \item \textbf{Consultas}, que necesitan respuesta inmediata: HTTP síncrono, con tiempo máximo de espera.
  \item \textbf{Avisos}, que toleran retraso, como las notificaciones: mensajería asíncrona con confirmación de recepción y cola de mensajes fallidos, de modo que un aviso no se pierda si el servicio receptor está caído.
\end{itemize}

Esta propuesta responde a la limitación del aviso \emph{pendiente} y se relaciona con el requisito RNF-04 (manejo de fallos). Está sin implementar: el esqueleto actual solo usa HTTP.
